\section{Experiments}
\label{sec:experiments}

\subsection{Dataset}

We use the \emph{GM12878} (Epstein--Barr virus-immortalised B-lymphocyte cell line)
dataset from the original DeepChrInteract benchmark~\cite{deepchrinteract_doc},
derived from publicly available Hi-C contact maps~\cite{hic}.
The dataset contains paired enhancer--promoter sequences of length $L = 10{,}001$~bp,
balanced between positive (interacting) and negative (non-interacting) pairs.
Full dataset statistics are reported in \Cref{tab:dataset}.

\begin{table}[H]
\centering
\caption{Dataset statistics for GM12878.}
\label{tab:dataset}
\begin{tabular}{lrrrr}
\toprule
Split & Positive & Negative & Total & Fraction \\
\midrule
Training   & --- & --- & --- & 80\% \\
Validation & --- & --- & --- & 10\% \\
Test       & --- & --- & --- & 10\% \\
\bottomrule
\end{tabular}
\end{table}

\textit{Note: statistics will be filled in upon execution of the preprocessing script on
the raw data.}

\subsection{Evaluation Metrics}

We evaluate all models on three primary metrics:

\begin{description}[noitemsep]
  \item[AUROC] Area under the receiver operating characteristic curve.
        The primary metric used in the original DeepChrInteract and SPEID papers.
  \item[AUPRC] Area under the precision--recall curve.
        More sensitive than AUROC for imbalanced datasets and preferred in
        recent EPI literature.
  \item[F1 Score] Harmonic mean of precision and recall at decision threshold 0.5.
\end{description}

Each configuration is trained five times with different random seeds
($\{0,1,2,3,4\}$) and we report mean $\pm$ standard deviation.

\subsection{Training Protocol}

\begin{itemize}[noitemsep]
  \item Optimiser: Adam~\cite{adam}, learning rate $5 \times 10^{-5}$, weight decay $10^{-5}$.
  \item Batch size: 32.
  \item Learning rate schedule: cosine annealing with warm restarts ($T_0 = 50$ epochs).
  \item Early stopping: patience 15 epochs, monitored on validation AUROC.
  \item Maximum epochs: 200.
  \item GPU: experiments run on a single NVIDIA GPU with at least 8~GB VRAM.
  \item For LLM fine-tuning: LLM layer learning rate $5 \times 10^{-6}$,
        head learning rate $5 \times 10^{-5}$ (layerwise learning rate decay).
\end{itemize}

\subsection{Ablation Study: Model Configurations}

\Cref{tab:ablation} lists the eight primary experimental configurations.

\begin{table}[H]
\centering
\caption{Experiment configurations.
         ``1B'' = single branch (E+P concatenated as input);
         ``2B'' = dual branch (independent encoders + fusion);
         fusion default = \textsc{Concat+Sub+Mul}.}
\label{tab:ablation}
\begin{tabular}{clllccc}
\toprule
ID & Encoder & Arch. & Branch & AUROC & AUPRC & F1 \\
\midrule
E01 & One-Hot  & CNN           & 1B & --- & --- & --- \\
E02 & One-Hot  & CNN           & 2B & --- & --- & --- \\
E03 & $k$-mer  & CNN           & 2B & --- & --- & --- \\
E04 & One-Hot  & BiLSTM        & 2B & --- & --- & --- \\
E05 & One-Hot  & Transformer   & 2B & --- & --- & --- \\
E06 & One-Hot  & CNN+BiLSTM    & 2B & --- & --- & --- \\
E07 & One-Hot  & CNN+Transformer & 2B & --- & --- & --- \\
E08 & DNABERT-2  & LLM (frozen) & 2B & --- & --- & --- \\
\midrule
E09 & HyenaDNA & LLM (fine-tune) & 2B & --- & --- & --- \\
\bottomrule
\end{tabular}
\end{table}

\subsection{Fusion Strategy Ablation}

Using E02 (CNN dual-branch) as the base model, we compare all six fusion strategies
from \Cref{tab:fusion} on the GM12878 test set:

\begin{table}[H]
\centering
\caption{Fusion strategy comparison (E02 base, GM12878 test set).}
\label{tab:fusion_abl}
\begin{tabular}{lccc}
\toprule
Fusion & AUROC & AUPRC & F1 \\
\midrule
Concat        & --- & --- & --- \\
Add           & --- & --- & --- \\
Subtract      & --- & --- & --- \\
Multiply      & --- & --- & --- \\
Bilinear      & --- & --- & --- \\
Concat+Sub+Mul & --- & --- & --- \\
\bottomrule
\end{tabular}
\end{table}

\subsection{Cross-Cell-Type Generalisation}

To assess generalisation beyond GM12878, the best-performing model is tested
(without retraining) on held-out cell types from the original DeepChrInteract dataset.
This evaluates whether the learned sequence grammar transfers across cell type contexts,
as discussed in~\cite{speid}.

\subsection{Expected Results and Baselines}

The original DeepChrInteract~\cite{deepchrinteract_doc} reported AUROC values in the
range of 0.70--0.83 on GM12878, with the CNN dual-branch (one-hot) as the strongest
non-LLM baseline.
SPEID~\cite{speid} reported AUROC of $\sim$0.78 on similar data using CNN+LSTM.
We anticipate that:
\begin{enumerate}[noitemsep]
  \item The BiLSTM and Transformer hybrid models (E06, E07) will outperform the
        pure CNN baseline (E02) by 1--3 AUROC points.
  \item The LLM-based encoders (E08, E09), particularly HyenaDNA (fine-tuned),
        will achieve the highest AUROC owing to their pre-training on large
        genomic corpora.
  \item The \textsc{Concat+Sub+Mul} fusion will consistently outperform simple
        concatenation by $\geq 0.5$ AUROC points.
\end{enumerate}
